%% kapa: exam=1 round=19 year=2008 period=1 form=A subject=appraisal_law subject_ko=부동산관계법규 seq=1 number=41
%% source: Q-Net 감정평가사 기출문제 (cst00309) artlSeq=1001962 file=extracted/1008/1001962/2008년 제19회 감정평가사 1차시험 기출문제/1교시.hwp
%% conversion: kordoc 4.24.51 -> tools/gen1.py (auto); review=proofread-sonnet; flags=none
\begin{question}{41}
지구단위계획구역 안에서의 건폐율 등의 완화적용에 관한 설명으로 옳지 않은 것은?
\begin{choices}
  \item 제1종지구단위계획구역 안에서 건축물을 건축하려는 자가 그 대지의 일부를 학교용지로 제공하는 경우에는 당해 건축물에 대하여 법령에서 정하는 비율까지 건폐율을 완화하여 적용할 수 있다.
  \item 제1종지구단위계획구역 안에서 도시계획조례의 규정에도 불구하고 법률에서 허용하는 건폐율의 범위 안에서 제1종지구단위계획으로 건폐율을 완화하여 적용할 수 있다.
  \item 제1종지구단위계획구역의 지정목적이 부설주차장을 설치하고자 하는 경우에는 주차장법에 의한 주차장 설치기준을 100\%까지 완화하여 적용할 수 있다.
  \item 계획관리지역에 지정된 제2종지구단위계획구역 안에서는 제2종지구단위계획으로 건축물의 용도․종류 및 규모 등을 완화하여 적용할 수 있다.
  \item 제2종지구단위계획구역 안에서는 제2종지구단위계획으로 당해 용도지역 또는 개발진흥지구에 적용되는 건폐율의 150\% 및 용적률의 200\% 이내에서 완화하여 적용할 수 있다.
\end{choices}
\end{question}
\answer{3}
